\section{Слой Минковского $(3{+}1)$ и дискретный носитель}
\label{sec:minkowski-layer}

Модель не строит пространство-время как четырёхмерную решётку $L^4$
и не выводит размерность $(3{+}1)$ компактификацией дополнительных измерений.
\textbf{Слой Минковского} --- тот же каркас, на котором уже записаны СИ и лабораторная физика:
три пространственные оси, одна временная, световой конус и скорость~$c$.
На этом слое дискретность вводится \emph{внутри} трёхмерного пространства, а не вместо него.

\begin{definition}[Разделение $(3{+}1)$ и носителя на~$M$]
\label{def:minkowski-carrier-split}
\begin{itemize}
  \item \textbf{Пространство} --- $\mathbb{R}^3$; носитель $\carrierLattice\subset\mathbb{R}^3$
  (\latGCK, \cref{thm:fcc-choice}).
  \item \textbf{Время} --- дискретные такты $t\in\mathbb{Z}$ с шагом~$\htick$; координата $t$ не четвёртая ось ячейки.
  \item \textbf{Каузальность} --- за один такт сигнал не дальше одного шага~$\caStepSpace$ по рёбрам $N(x)$
  (аксиома~$A_1$, \cref{ax:A1}); микроскопический конус задаётся $\caSignalSpeed=\caStepSpace/\htick$.
  \item \textbf{Макроскопический $c$} --- осреднённый null на~$\layerMacro$,
  $\speedC=\caKappaGeom \caSignalSpeed$ с якорем на полном $(3{+}1)$-носителе (\cref{sec:two-c,sec:c-light-cone}).
\end{itemize}
\end{definition}

Отказ от непрерывного предела $\caStepSpace,\htick\to 0$ на~$\layerMicro$
(\cref{sec:no-continuum-carrier}) не отменяет слой Минковского:
он запрещает трактовать микромир как гладкое $\mathbb{R}^4$, а не отбрасывает $(3{+}1)$ как язык прибора.
Непрерывные поля и дифференциальные уравнения остаются языком~$\layerMacro$ на \emph{фиксированной} решётке
с тем же конусом и теми же размерностями SI.

Гексагональный срез $(2{+}1)$ (\cref{sec:hex-slice,sec:31-21-slice})
--- не альтернативная физика, а срез того же макроскопического $(3{+}1)$-bulk:
мировые линии и поле могут быть ограничены плоскостью плотной упаковки,
но якорь $\speedC$ и $|N|=12$ задаются полной \latGCKshort\ в объёме.

\medskip
\noindent
Итог для читателя: \emph{стандартный} Минковский $(3{+}1)$ принят по построению;
задача модели --- показать, какой обратимый закон~$g$ на $\carrierLattice\times\mathbb{Z}$
воспроизводит на~$\layerMacro$ уже сверенную физику, а не «вырастить» пространство из абстрактного графа.
